%hw-1.tex
%%%%%%%%%%%%%%%%%%%%%%%%%%%%%%%%%%%%%%%%%%%%%%%%%%%%%%%%%%%%%%%%%%%%%%
%%%%%%            Math 403 Homework \#1, Spring 2025            %%%%%%
%%%%%%                                                          %%%%%%
%%%%%%                 Instructor: Ezra Miller                  %%%%%%
%%%%%%%%%%%%%%%%%%%%%%%%%%%%%%%%%%%%%%%%%%%%%%%%%%%%%%%%%%%%%%%%%%%%%%
\documentclass[11pt]{article}


\oddsidemargin=17pt \evensidemargin=17pt
\headheight=9pt     \topmargin=26pt
\textheight=564pt   \textwidth=433.8pt
%voffset=-6ex

\usepackage{amsmath,amssymb,graphicx,color,enumitem}%,setspace,mathrsfs}%,hyperref
\setlist[enumerate]{parsep=1.5ex plus 4pt,topsep=-1ex plus 4pt,itemsep=0ex}
  \setlist[itemize]{parsep=1.5ex plus 4pt,topsep=-1ex plus 4pt,itemsep=-1.33ex}

\leftmargini=5.5ex
\leftmarginii=3.5ex

%for \marginpar to fit optimally
%hoffset=-1.02in
\setlength\marginparwidth{2.2in}
\setlength\marginparsep{1mm}
\newcommand\red[1]{\marginpar{\linespread{.85}\sf%
		   \vspace{-1.4ex}\footnotesize{\color{red}#1}}}
\newcommand\score[1]{\marginpar{\vspace{-2ex}\color{blue}{#1/3}}\hspace{-1ex}}
\newcommand\extra[1]{\marginpar{\color{blue}{#1/1}}\hspace{-1ex}}
\newcommand\total[2]{\marginpar{\colorbox{yellow}{\huge #1/#2}}}
\newcommand\collab[1]{\marginpar{\vspace{-11ex}\colorbox{yellow}{#1/3}}\hspace{-1ex}}
\newcommand\magenta[1]{\colorbox{magenta}{\!\!#1\!\!}}
\newcommand\yellow[1]{\colorbox{yellow}{\!\!#1\!\!}}
\newcommand\green[1]{\colorbox{green}{\!\!#1\!\!}}
\newcommand\cyan[1]{\colorbox{cyan}{\!\!#1\!\!}}
\newcommand\rmagenta[2][0ex]{\red{\vspace{#1}\magenta{\phantom{:}}\,: #2}}
\newcommand\ryellow[2][0ex]{\red{\vspace{#1}\yellow{\phantom{:}}\,: #2}}
\newcommand\rgreen[2][0ex]{\red{\vspace{#1}\green{\phantom{:}}\,: #2}}
\newcommand\rcyan[2][0ex]{\red{\vspace{#1}\cyan{\phantom{:}}\,: #2}}

%For separated lists with consecutive numbering
\newcounter{separated}

\newcommand{\excise}[1]{}
\newcommand{\comment}[1]{{$\star$\sf\textbf{#1}$\star$}}

%%%%%%%%%%%%%%%%%%%%%%%%%%%%%%%%%%%%%%%%%%%%%%%%%%%
%                                                 %
% TO ENABLE GRADING, DO NOT ALTER ABOVE THIS LINE %
%                                                 %
%%%%%%%%%%%%%%%%%%%%%%%%%%%%%%%%%%%%%%%%%%%%%%%%%%%

%new math symbols taking no arguments
\usepackage{quiver}
\usepackage{commath}
\usepackage{amsmath}
\newcommand\0{\mathbf{0}}
\newcommand\CC{\mathbb{C}}
\newcommand\NN{\mathbb{N}}
\newcommand\QQ{\mathbb{Q}}
\newcommand\RR{\mathbb{R}}
\newcommand\ZZ{\mathbb{Z}}
\newcommand\bb{\mathbf{b}}
\newcommand\kk{\Bbbk}
\newcommand\mm{\mathfrak{m}}
\newcommand\xx{\mathbf{x}}
\newcommand\yy{\mathbf{y}}
\newcommand\minus{\smallsetminus}
\newcommand\goesto{\rightsquigarrow}

%redefined math symbols taking no arguments
\newcommand\<{\langle}
\renewcommand\>{\rangle}
\renewcommand\iff{\Leftrightarrow}
\renewcommand\implies{\Rightarrow}

%to explain about LaTeX commands:
\newcommand\command[1]{\texttt{$\backslash$#1}}

%new math symbols taking arguments
\newcommand\ol[1]{{\overline{#1}}}

%redefined math symbols taking arguments
\renewcommand\mod[1]{\ (\mathrm{mod}\ #1)}

%roman font math operators
\DeclareMathOperator\im{im}

%for easy 2 x 2 matrices
\newcommand\twobytwo[1]{\left[\begin{array}{@{}cc@{}}#1\end{array}\right]}

%for easy column vectors of size 2
\newcommand\tworow[1]{\left[\begin{array}{@{}c@{}}#1\end{array}\right]}


%%%%%%%%%%%%%%%%%%%%%%%%%%%%%%%%%%%%%%%%%%%%%%%%%%%%%%%%%%%%%%%%%%%%%%
\begin{document}%%%%%%%%%%%%%%%%%%%%%%%%%%%%%%%%%%%%%%%%%%%%%%%%%%%%%%
%%%%%%%%%%%%%%%%%%%%%%%%%%%%%%%%%%%%%%%%%%%%%%%%%%%%%%%%%%%%%%%%%%%%%%


\title{\mbox{}\\[-8ex]Math 403 Homework \#1, Spring 2025\\\normalsize
	Instructor: Ezra Miller\\[-2.5ex]}
\author{Solutions by: Thomas Kidane\\[1ex]
Collaborators: Tymon Hendershott, Holly Keegan, Ben Greene, Ricardo Prado Cunha\\
(1 point for each of up to 3 collaborators who also list you)}
\date{Due: 11:59pm Saturday 25 January 2025}

\maketitle

\vspace{-5ex}\collab{}%

\noindent
\textsc{Reading assignments}\total{}{57}% cover Lectures 1 - 7
\hfill
(item numbers are lecture numbers as in PDF lecture notes)
\begin{enumerate}[itemsep=-1ex]
\item% 1.
for Thu.~9 January
  \begin{itemize}
  \item{}%
  [Hefferon, p.153--154] field axioms
  \item{}%
  [Cornell, ``Fields'' (about 1/3 of the way through
  \texttt{4330-week1.pdf})]: more on fields%
  \item{}%
  [Climenhaga, \S 7.2, \S 9.1] isomorphism, rank-nullity
  \end{itemize}

\item% 2.
for Tue.~14 January
  \begin{itemize}
  \item{}%
  [Lax, Chapter~1] quotients
  \item{}%
  [Climenhaga, \S 5.1] quotients
  \item{}%
  [Cornell, ``Quotient Spaces'' (\texttt{4330-week5.pdf})] universal properties
  \item{}%
  [Cornell, ``Exact Sequences'' (about halfway through \texttt{4330-week5.pdf})]
  \end{itemize}

\item% 3.
for Thu.~16 January
  \begin{itemize}
  \item{}%
  [Lax, Chapter~2] duality
  \item{}%
  [Climenhaga, \S 5.1, \S 8.2] duality, transpose
  \item{}%
  [Treil, \S 5.1--5.5] Hermitian inner products, ajoints; \S
  5.2--5.4 should be mostly review
  \end{itemize}

\item% 4.
for Tue.~21 January
  \begin{itemize}
  \item{}%
  [Lax, Appendix 15] Jordan canonical form (very short proof)
  \item{}%
  [Hefferon, \S 5.IV.1] characteristic polynomial and minimal polynomial
  \item{}%
  [Treil, \S 9.3 and \S 9.5] generalized eigenspaces and Jordan canonical form
  \end{itemize}

\item% 5.
for Thu.~23 January
  \begin{itemize}
  \item{}%
  [Lax, Chapter 14, p.214--221] norms, equivalence, continuity, local compactness
  \item{}%
  [Stewart--Sun, \S II.1] norms, equivalence, Hahn--Banach theorem
  \end{itemize}

\item% 6.
for Tue.~28 January
  \begin{itemize}
  \item{}%
  [Lax, Chap.12, p.187--190] Convex sets
  \end{itemize}

\item% 7.
for Thu.~30 January
  \begin{itemize}
  \item{}%
  [Serge Lang, \emph{Linear Algebra}, Chapter XII, \S1--\S2]
  separating hyperplanes
  \item{}%
  [Serge Lang, \emph{Linear Algebra}, Chapter XII, \S3--\S4]
  support hyperplane, extreme~point$\!\!$
  \end{itemize}
\end{enumerate}
  
\pagebreak

\noindent
\textsc{Exercises}% cover Lectures 1 - 4

\begin{enumerate}
\item\ \score{}%1.
(Freshman's Dream): The \emph{characteristic} of a field~$F$ is the
smallest positive integer $p$ such that the sum $1 + \cdots + 1$ of
$p$ multiplicative identities is $0$ in~$F$.  Prove that $p$ is prime
if it is finite.  If $F$ has characteristic~$p$, show that $(a + b)^p
= a^p + b^p$ for $a,b \in F$.

\textbf{Solution}
\linebreak
For the sake of notional ease take the symbol 2 to mean 1+1.
\begin{align*} 
    2:&=1+1 \\ 
    3:&=1+1+1 \\
    4:&=1+1+1+1\\
    \dots\\
    p:&=0=1+1+1\dots
\end{align*}

From now on all the operations, identities, and all the inverses are from the field $F$.

We first see that by using the axioms of a field, multiplication with this notation is like modulo multiplication.

Armed with our notation, let's attack the question.

Assume $F$ is finite and its characteristic is a composite then there are two integeres $a,b$ such that $p = a*b$. Thne you
can find two elements $a^{\prime}, b^{\prime}$ such they are 1 just added $a$ and $b$ times respectively.


But this also implies that 0 would have an inverse because $a^{\prime} * b^{\prime}* b^{\prime-1}=a^{\prime}=0*b^{\prime-1}=0$, so this implies 0=$b^{\prime-1}$. But this is obviously nonsense.

So if $F$ is finite $p$ has to be prime.

Going to the second question, we can use the binomial theorem to get that
\begin{align}
    (a'+b')^p=\sum_{k=0}^{p} {  \binom{p}{k}}(a')^{p}(b')^{p-k}
\end{align}
If we look at the terms with $k\neq 0$ and $k \neq p$, we get that ${p \choose k}= p*((p-1)!/(p-k)!k!)$, since $p$=0, then all these terms evalute to 0 giving $(a'+b')^{p}=(a')^p+(b')^p$

\pagebreak


\item\ \score{}%2.
Fix a vector space $V$ and a subspace $W \subseteq V$ over a
field~$F$.  Let $\pi: V \to V/W$ be the \emph{projection homomorphism}
given by $\pi(v) = v + W$.  Write $X$ for the set of all subspaces
of~$V$ that contain~$W$, and write $Y$ for the set of all subspaces
of~$V/W$.  Prove that $\pi$ induces a bijection between these two
sets, with
\begin{align*}
X &\to Y
\\
L &\mapsto \pi(L) = \{\pi(v) \mid v \in L\}\\
\intertext{and}\\[-8.75ex]
Y &\to X
\\
M &\mapsto \pi^{-1}(M) = \{v \in V \mid \pi(v) \in M\}.
\end{align*}
\textbf{Solution} \\\\
Let \textit{L} $\in X$, then $W$ $\subseteq$ $L$.

Then $dim(L)\geq dim(W)$. So assume $dim(L)$ = $x$ and $dim(W)=y$, $\Rightarrow x\geq y$.
$\to$ $v_1, ... , v_x $ span $L$.

Change this vectors so that $v_1, ... , v_y $ span $W$ and the $y-x$ other vectors span the complement of $W$ in $L$. Now if you take the subspace spanned by $[v_{y+1}],[v_{y+2}],[v_{y+3}],...,[v_{x}]$, they are a subspace of $V/W$, and this subspace and $L$ are the equivalent(upto isomorphism), because the basis of this space is the basis for $W$ plus the other $x-y$ vectors, which is the same as $L$.

Take $M\in Y$, then it has a basis of say $x$ vectors. If you take the union of the elements in the basis(one elment from the coset) and the basis for $W$, you get an element in $Y$. So $\pi^{-1}$ is surjective. \\ \\
We have proved that $\pi$ and $\pi^{-1}$ are surjective, but we also need to prove that they are also injective.

For $\pi$, assume $L_1$ and $L_2$ both map to $M\in Y$. Now look at the basis for $M$, they are the basis for $W$ plus $v_1$ up to $v_y$ assuming $dim(M)=y$. Since $M=L_1$ and $M=L_2$, then $L_1=L_2$ because they have the same set of basis, therefore are the same subspace.

For $\pi^{-1}$, assume $M_1$ and $M_2$, both map to $L\in Y$. Since $M_1=L$ and $M_2=L$, then $M_1=M_2$, by the same reasoning as above. Henceforth, the function is bijective.

\newpage

\item\ \score{}%3.
Show that giving an exact sequence $\cdots \to V_{i-1} \to V_i \to
V_{i+1} \to \cdots$ is the same as giving a collection of short exact
sequences $0 \to K_i \to V_i \to K_{i+1} \to 0$, one for each~$i$.
(The long exact sequence is said to be constructed by \emph{splicing}
the short exact sequences together.)

For the sake of simplicity, I will use the diagram below.

% https://q.uiver.app/#q=WzAsMTQsWzAsMCwiMCJdLFsxLDEsIktfaSJdLFsyLDIsIlZfaSJdLFszLDMsIktfe2krMX0iXSxbNCw0LCIwIl0sWzIsNCwiMCJdLFs0LDIsIlZfe2krMX0iXSxbNSwxLCJLX3tpKzJ9Il0sWzQsMCwiMCJdLFs2LDAsIjAiXSxbNiwyLCJWX3tpKzJ9Il0sWzIsMCwiMCJdLFswLDIsIi4uLiJdLFs3LDMsIi4uLiJdLFswLDFdLFsyLDMsIiIsMCx7InN0eWxlIjp7ImhlYWQiOnsibmFtZSI6ImVwaSJ9fX1dLFszLDRdLFs1LDNdLFszLDYsIiIsMCx7InN0eWxlIjp7InRhaWwiOnsibmFtZSI6Imhvb2siLCJzaWRlIjoidG9wIn19fV0sWzEsMiwiIiwwLHsic3R5bGUiOnsidGFpbCI6eyJuYW1lIjoiaG9vayIsInNpZGUiOiJ0b3AifX19XSxbNiw3LCIiLDAseyJzdHlsZSI6eyJoZWFkIjp7Im5hbWUiOiJlcGkifX19XSxbNyw4XSxbNyw5XSxbNywxMCwiIiwwLHsic3R5bGUiOnsidGFpbCI6eyJuYW1lIjoiaG9vayIsInNpZGUiOiJ0b3AifX19XSxbMSwxMV0sWzEyLDFdLFsxMCwxMywiIiwwLHsic3R5bGUiOnsiaGVhZCI6eyJuYW1lIjoiZXBpIn19fV1d
\[\begin{tikzcd}
	0 && 0 && 0 && 0 \\
	& {K_i} &&&& {K_{i+2}} \\
	{...} && {V_i} && {V_{i+1}} && {V_{i+2}} \\
	&&& {K_{i+1}} &&&& {...} \\
	&& 0 && 0
	\arrow[from=1-1, to=2-2]
	\arrow[from=2-2, to=1-3]
	\arrow[hook, from=2-2, to=3-3]
	\arrow[from=1-5, to=2-6]
	\arrow[from=2-6, to=1-7]
	\arrow[hook, from=2-6, to=3-7]
	\arrow[from=3-1, to=2-2]
	\arrow[two heads, from=3-3, to=4-4]
	\arrow[two heads, from=3-5, to=2-6]
	\arrow[two heads, from=3-7, to=4-8]
	\arrow[hook, from=4-4, to=3-5]
	\arrow[from=4-4, to=5-5]
	\arrow[from=5-3, to=4-4]
\end{tikzcd}\]


I highlighted that the mapping between $K_i$ and $V_i$ is injective and that the mapping between $V_i$ and $K_{i+1}$ is surjective. Now consider what being an exact sequence entails, it just means that the you can make a mapping satisfying $img(\varphi_i)=ker(\varphi_{i+1})$. So if we can construct a mapping between the vector spaces that satisfies this property, then the aforementioned exact sequence does in fact exist.

First step let's make $\varphi_i$, such that its domain is $V_i$ and its codomain is $V_{i+1}$. We can do this by first observing that there exists a surjective mapping between $V_i$ and $K_i$. Next we can also observe that there is an injective mapping between $K_{i}$ and $V_{i+1}$. If we compose these to mappings, we get a mapping whose domain is $V_i$ and its codomain is $V_{i+1}$.

Next let's make $\varphi_{i+1}$ whose domain is $V_{i+1}$ and whose codomain is $V_{i+2}$, such that its kernel is the image of the previous mapping. We know that from the smaller chain consisting of $V_{i+1}$ that there is a mapping from $V_{i+1}$ to $K_{i+2}$, such that it sends all the elements that have a preimage in $K_{i+1}$ to \textbf{0}. We also now that there is an injective mapping between $K_{i+2}$ and $V_{i+2}$. If we compose this two mappings, we get a mapping whose domain is $V_{i+1}$ and whose codomain is $V_{i+2}$, such that its kernel is the image of $K_{i+1}$ in $V_{i+1}$.

So now we have that the $img(\varphi_{i}) \subseteq ker(\varphi_{i+1})$, let's now prove that every element in the kernel of $\varphi_{i+1}$ has a preimage in $V_i$, thereby completing the proof.

We know from the chain consisting of $V_{i+1}$, that the kernel of the mapping between $V_{i+1}$ and $K_{i+2}$, is $K_{i+1}$ by definition of exactness. We also now that each element in $K_{i+1}$ has a preimage in $V_i$, so to conclude the kernel of the mapping between $K_{i+2}$ and $V_{i+2}$ is $K_{i+1}$ and every element in $K_{i+1}$ has a preimage in $V_i$. So $ker(\varphi_{i+1}) = img(\varphi_i)$. So given a collection of short sequences, you
can make a longe chain. To do the reverse, take each the kernel of the homomorphism beteween $V_{i}$ and $V_{i+1}$ and define it as $K_{i+1}$. Now if you constrain
homomorphism to this codomain, you get the short sequences again. Therefore a collection of short sequences is the same as one big long sequence.

\newpage
\item\ \score{}\label{rn}%4.
Rank-nullity theorem for exact sequences: Given an exact sequence
$$%$$ $
  0 \to V_0 \to V_1 \to \cdots \to V_r \to 0,
$$%$$ $
prove that $\sum_{i=0}^r (-1)^i \dim V_i = 0$.

\textbf{Solution}

For the sake of simplicity I will extend the exact sequence on the right by adding another \textbf{0}, and I will relabel the sequences as shown below.

\[\begin{tikzcd}
	0 & {V_0} & {V_1} & {...} & {V_r} & 0 & 0
	\arrow["{\varphi_0}", from=1-1, to=1-2]
	\arrow["{\varphi_1}", from=1-2, to=1-3]
	\arrow["{\varphi_2}", from=1-3, to=1-4]
	\arrow["{\varphi_{r}}", from=1-4, to=1-5]
	\arrow["{\varphi_{r+1}}", from=1-5, to=1-6]
	\arrow["{\varphi_{r+2}}", from=1-6, to=1-7]
\end{tikzcd}\]

\begin{align*}
    dim(V_i)&=img(\varphi_{i+1}) + ker(\varphi_{i+1})\\
    dim(V_i)&=ker(\varphi_{i+2})+ker(\varphi_{i+1})\\ \\
    dim(V_0)&=ker(\varphi_2)+ker(\varphi_1)\\
    (-1)dim(V_1)&=(-1)(ker(\varphi_3)+ker(\varphi_2))\\
    ...\\
    (-1)^r dim(V_r)&=(-1)^r(ker(\varphi_{r+2})+ker(\varphi_{r+1}))\\
    \sum_{i}^{r}dim(V_i)&=ker(\varphi_1)+(-1)^{r}ker(\varphi_{r+2})
\end{align*}

But $ker(\varphi_1)=0$ and $ker(\varphi_{r+2})=0$, so their sum is 0, thereby completing the proof.

\pagebreak

\item\ \score{}%5.
Rank-nullity for arbitrary complexes: Given a complex $0 \to V_0 \to
V_1 \to \cdots \to V_r \to 0$, write $B_i = \im(V_{i-1} \to V_i)$ and
$Z_i = \ker(V_i \to V_{i+1})$.  Prove that
$$%$$ $
  \sum_{i=0}^r (-1)^i \dim V_i = \sum_i (-1)^i \dim H_i,
$$%$$ $
where $H_i = Z_i/B_i$ is the $i^\text{th}$ homology of the complex.
[In Exercise~\ref{rn}, $B_i = Z_i = K_i$, so $H_i = Z_i/B_i = 0$ for
all~$i$.]

\textbf{Solution}

From corollary of lecutre 2, $dim(H_i)=dim(Z_i)-dim(B_i)$.

\begin{align*}
  \sum_{i=0}^{r}(-1)^{i}V_i&=dim(Z_0)+dim(B_1)\\
  &-dim(Z_1)-dim(B_2)\\
  \dots\\
  &+(-1)^{r}(dim(Z_r)+dim(B_{r+1}))
\end{align*}

\begin{align*}
  \sum_{i=0}^{r}(-1)^{i}H_i&=dim(Z_0)-dim(B_0)\\
  &-dim(Z_1)+dim(B_1)\\
  \dots\\
  &+(-1)^{r}(dim(Z_r)-dim(B_{r}))
\end{align*}

\begin{align*}
  \sum_{i=0}^{r}(-1)^{i}V_i-\sum_{i=0}^{r}(-1)^{i}H_i &= (dim(Z_0)-dim(Z_0))
   +(dim(B_1)-dim(B_1)) + dim(B_0) \\ &+ (-dim(Z_1)+dim(Z_1)) + \dots + dim(B_{r+1})\\
   &=dim(B_0)+dim(B_{r+1})=0+0=0
\end{align*}

So the sums are equal, because $dim(B_0)$ is 0 because 0 can only be mapped to 0 in $V_1$ and 
$dim(B_r)$=0 because $V_r$ is mapped to 0, so their sum is 0.

\pagebreak

\item\ \score{}%6.
Two elements $u$ and~$v$ in a vector space~$V$ are \emph{congruent
modulo} a subspace~$W \subseteq V$, written $u \equiv v \mod W$, if $u
+ W = v + W$.  Show that congruence modulo~$W$ is an \emph{equivalence
relation} on~$V$, meaning that it is
\begin{itemize}
\item
reflexive: $v \equiv v \mod W$ for all $v \in V$;
\item
symmetric: if $u \equiv v \mod W$, then $v \equiv u \mod W$; and
\item
transitive: if $u \equiv v \mod W$ and $v \equiv x \mod W$ then $u
\equiv x \mod W$.
\end{itemize}

\textbf{Solution}
\\
Reflexivity: $v+0=v=v$ and $0\in W$, so $v\equiv v$(mod $W$). \\
Symmetry: $u\equiv v$(mod $W$)$\rightarrow \exists w$, $w\in W$ such that $u+w=v$ $\rightarrow$ $u=v-w \rightarrow v-w=u \rightarrow v\equiv u$(mod $W$) because $-w \in W$\\
Transitivity: $u\equiv v$(mod $W$) $\rightarrow \exists w$, such that $u+w_1=v$. Same for $v\equiv x$, there $\exists w_2 \in W$ such that $v+w_2=x$. And since $w_1+w_2 \in W$($W$ is closed under addition), $u+w_1+w_2=v+w_2=x$, so $u\equiv x$(mod $W$).

\pagebreak

\item\ \score{}%7.
Give an example of three subspaces $Y_1$, $Y_2$, and $Y_3$ in $\RR^2$
such that $Y_1 + Y_2 + Y_3 = \RR^2$ and $Y_i \cap Y_j = \{\0\}$ for
all $i \neq j$, but $\RR^2$ is not the direct sum of $Y_1$, $Y_2$,
and~$Y_3$.

\textbf{Solution}

Take $Y_1=span((0,1)) \quad Y_2=span((1,0)) \quad Y_3=span((1,1))$
The pairwise intersection is $\{0\}$, their sum is $\RR^2$ but their direct sum isn't $\RR^2$, because their intersection isn't trivial.

\pagebreak

% 
\item\ \score{}%8.
Prove that if $V$ is a vector space and $W \subseteq V$ is a subspace,
then $W$ has a \emph{complement}: a subspace $U \subseteq V$ such that
$V = W \oplus U$.  Hint: $V/W$ has a basis; lift it back to $V$.

\textbf{Solution}
Take the basis for $V/W$. Take one element(non-zero) from each basis(coset) and the union of these with the set of the basis for 
$W$ will be a basis for $V$.

Proof: Everytime you take an elmenent from the cosets, apply gram-schidt(Assume you started with an orthogonal set of basis for $W$) process to it, and then you will end
up with an orthogonal set of basis. This is because you will never get a zero during this process,
if you got zero during this process, that would imply that the vector you used is a linear combitionaion 
of the other ones and therefore couldn't be a basis for $V/W$.

After you have made this set of basis the elements that are not a basis for W, will be a a basis for $U$. Their direct sum will also
be $V$, since they have no intersection.


\pagebreak


\item%9.
For vectors $\xx = (1, 2i, 1+i)$ and $\yy = (i, 2-i, 3)$, compute
\begin{enumerate}[itemsep=-1ex]
\item\ \score{}%
$\<\xx,\yy\>$, $||\xx||^2$, $||\yy||^2$, and $||\yy||$;
\item\ \score{}%
$\<3\xx,2i\yy\>$ and $\<2\xx,i\xx + 2\yy\>$;
\item\ \score{}%
$||\xx+2\yy||$.  [Use parts (a) and~(b) for this.]
\end{enumerate}

\textbf{Solution}

\begin{enumerate}
  \item 
  $\<\xx,\yy\> = \bar y^{T}x = (-i,2+i,3)^{T} \cdot (1,2i,1+i) = -i + 4i-2+3+3i=1+6i$. 
  \\ $\norm{\vec{x}}^2=\bar x^{T}x= (1,-2i,1-i)^T \cdot (1,2i,1+i) = 1+4+1-i+i+1=7$\\
  $\norm{\vec{y}}^2=\bar y^{T}y= (-i,2+i,3)^T \cdot (i,2-i,3) =1+4+1+9 =15$
  \\
  $||y|| = \sqrt{15}$
  \item 
  $\<3\xx,2i\yy\> = -6i\<\xx,\yy\>=-6i+36$\\
  $\<2\xx,i\xx + 2\yy\> = 2 \cdot (-i\<\xx,\xx\>+2\<\xx,\yy\>)=2 \cdot (-7i + 2+12i)=10i+4$

  \item 
  \begin{align*} 
    ||\mathbf{x} + 2 \mathbf{y} || & = \sqrt{\langle x + 2y, x + 2y \rangle} \\
                                   & = \sqrt{\langle x, x \rangle + \langle x, 2y \rangle + \langle 2y, x \rangle + \langle 2y, 2y \rangle}\\ 
                                   & = \sqrt{||x||^2 + ||2y||^2 + 2 \langle x, y \rangle + 2 \langle y, x \rangle } \\
                                   & = \sqrt{7+ 60 + 4}\\
                                   & = \sqrt{71}
  \end{align*} 

  
   \pagebreak
\end{enumerate}

\item\ \score{}%10.
Prove that for vectors in an inner product space,
$||\xx\pm\yy||^2 = ||\xx||^2 + ||\yy||^2 \pm 2\mathrm{Re}\<\xx,\yy\>$.

\textbf{Solution}

$\<\xx \pm \yy, \xx \pm \yy\> = \<\xx,\xx \pm \yy\> \pm \<\yy,\xx \pm \yy\> = (\<\xx,\xx\> \pm \<\xx,\yy\>) \pm (\<\yy,\xx\> \pm \<\yy,\yy\>)\\\ $
$||\xx||^2 + ||\yy||^2 \pm \<\xx,\yy\> \pm \overline{\<\xx,\yy\>} = ||\xx||^2 + ||\yy||^2 \pm 2Re(\<\xx,\yy\>)$

\pagebreak



\item\ \score{}%11.
For any $m \times n$ complex matrix~$A$, prove that $\ker(A^* A) =
\ker(A)$.

\textbf{Solution}

Let $\xx \in F^{n}$ such that  $\xx \in ker(A) \rightarrow A^{*}A\xx =  A^{*}(0) = 0 \rightarrow \xx \in ker(A^{*}A)$\\ 
Let $\xx \in F^{n}$ such that $A^{*}Ax=0 \rightarrow (\xx^{*}A^{*})A\xx = \<A\xx,A\xx\>  = 0 \rightarrow A\xx  = 0 \rightarrow \xx \in ker(A)$.

\pagebreak

\item\ \score{}%12.
Prove that if $P$ is self-ajoint (that is, $P^* = P$) and idempotent
(that is, $P^2 = P$) then $P$ is the matrix for an orthogonal
projection.

 \textbf{Solution}
  Defintion(I was not given any defintion) assumed:- P is orthogonal projection iff $P\xx  \xx \perp P\xx$\\
 $\<P\xx-\xx, P\xx\> = \<P^{*}(P\xx-\xx),\xx\> = \<(P^{2}\xx-P\xx),\xx\> = \<P\xx-P\xx,\xx\>=\<0,\xx\> = 0$
 Therefore P is orthogonal projection.
 
\pagebreak

\item\ \score{}%13.
If $V$ is a vector space over~$\CC$ of dimension~$n$, then it is also
a vector space over~$\RR$ of dimension~$2n$.  (If this isn't clear to
you, then write down a proof.)  Given a Hermitian inner product
$\<\xx,\yy\>$ on~$V$ as a complex vector space, show that the real
part $\<\xx,\yy\>_\RR = \mathrm{Re}(\<\xx,\yy\>)$ is an inner product
on~$V$ as a real vector space.



\textbf{Solution}

We can basically ask the question, if you only take the real part of the complex Hermitian inner product
will it give you a valid inner product for the real vector space. So we just need to check 
if it satisfies the properties of linearity, conjugate symmetry, nonegativity and nondegeneracy.

1. Linearity \\
This property is simply true because the linearity is satisfied in complex inner product, and if 
you take the real part of it, it also has to be linear otherwise it wouldn't be linear in the compolex inner product to begin with, or you 
could make a map from from the norm to the real numbers, and since the norm is linear and this function is linear, 
their composition is also linear.

2.Conjugate Symmetry \\
Just like Linearity, when you only take the real part, there has to be symmetry otherwise it wouldn't 
be symmetrical in the complex inner product anyway.

3.Nonnegativity \\
The product is already real in the complex inner product.

4.Nondegeneracy \\
The product is already real and nondegenerate in the complex inner product.

\pagebreak

\item\ \score{}%14.
Find all possible Jordan forms of linear transformations with
characteristic polynomial $(t - 1)^2 (t + 2)^2$.

\textbf{Solution}

$\begin{bmatrix}
  1 & 0 & 0 & 0\\
  0 & 1 & 0 & 0\\
  0 & 0 & -2 & 0\\
  0 & 0 & 0 & -2  
\end{bmatrix}$
$\begin{bmatrix}
  1 & 1 & 0 & 0\\
  0 & 1 & 0 & 0\\
  0 & 0 & -2 & 0\\
  0 & 0 & 0 & -2  
\end{bmatrix}$
$\begin{bmatrix}
  1 & 0 & 0 & 0\\
  0 & 1 & 0 & 0\\
  0 & 0 & -2 & 1\\
  0 & 0 & 0 & -2  
\end{bmatrix}$
$\begin{bmatrix}
  1 & 1 & 0 & 0\\
  0 & 1 & 0 & 0\\
  0 & 0 & -2 & 1\\
  0 & 0 & 0 & -2  
\end{bmatrix}$

Including permutaions on Jordan blocks.

\pagebreak

\item\ \score{}%15.
Find all possible Jordan forms of linear transformations with
characteristic polynomial $(t - 2)^3 (t + 1)$ and minimal polynomial
$(t - 2)^2 (t + 1)$.

$\begin{bmatrix}
  2 & 1 & 0 & 0\\
  0 & 2 & 0 & 0\\
  0 & 0 & 2 & 0\\
  0 & 0 & 0 & -1  
\end{bmatrix}$
Up to permutation of the Jordan blocks.
\pagebreak

\item\ \score{}%16.
How many similarity classes are there for $3 \times 3$ matrices whose
only eigenvalues are $-3$ and~$4$?

\textbf{Solution}
The answer is 4. There are at least 2 eigenvectors, with eigenvalues
-3 and 4. The third one can be in either and the geomtric multiplicies of the 
two eigenvalues can be either 1 or 2 for each, giving a total of 4 equivalence classes.

\pagebreak

\item\ \score{}%17.
Prove or disprove: two $n \times n$ matrices are similar if and only
if they have the same characteristic polynomial and minimal
polynomial.

\textbf{Solution}

Unfortunately, two matrix having the same characteristic and minimal polynomial,
doesn't gurantee that they are similar, it is a requirement but it is not a suffiecient requirement.

As an example consider a matrix whose eigenvalue has geomtric multiplicity 3, algebraic multiplicty 8(I didn't write the matrix out because it was 128 numbers),
so that you can get two matrices, such that one has 3 jordan blocks of size 4,2,2 and the other has 3 jordan blocks
of size 4,3,1. They have the same minimal polynomial and characteristic polynomial, but their jordan blocks are different
and therefore can't be the similar.

\end{enumerate}


%%%%%%%%%%%%%%%%%%%%%%%%%%%%%%%%%%%%%%%%%%%%%%%%%%%%%%%%%%%%%%%%%%%%%%
\end{document}%%%%%%%%%%%%%%%%%%%%%%%%%%%%%%%%%%%%%%%%%%%%%%%%%%%%%%%%
%%%%%%%%%%%%%%%%%%%%%%%%%%%%%%%%%%%%%%%%%%%%%%%%%%%%%%%%%%%%%%%%%%%%%%


%%% Various Emacs customizations:
%%% Local Variables:
%%% fill-column: 70
%%% indent-tabs-mode: t
%%% TeX-electric-sub-and-superscript: nil
%%% TeX-brace-indent-level: 0
%%% LaTeX-indent-level: 0
%%% LaTeX-item-indent: 0
%%% TeX-master: t
%%% End: